%%%%%%%%%%%%%%%%%%%%%%%%%%%%%%%%%%%%%%%%%%%%%%%%%%%%%%%%%%%%%%%%%%
%% Toto je inkludovaný LaTeXový soubor `ca_docop.tex'.          %%
%% ------------------------------------------------------------ %%
%% Verze 2022-05-31                                             %%
%% Vyvíjí & udržuje <stanislav.hledik@physics.slu.cz>           %%
%%%%%%%%%%%%%%%%%%%%%%%%%%%%%%%%%%%%%%%%%%%%%%%%%%%%%%%%%%%%%%%%%%

\chapter{Volby dokumentní třídy}\label{options}%%%%%%%%%%%%%%%%%%%%%%%%%%%%

Dokumentní třída\INDEX{dokumentní třída} \pgm{\ipotname .cls}\INDEX{dokumentní
  třída!ipothes} akceptuje následující ortogonální\footnote{Aktivace určité
  volby neovlivňuje funkčnost voleb ostatních (kromě \opt{draft} a
  \opt{final}).} volby:

\begin{description}
\optitem{draft}\label{opdraft} se používá pro často opakované průběžné
  kompilace během psaní práce. Přetečené řádky budou na pravé straně označeny
  černým obdélníčkem (\TeX isty nazývaným \emph{slimák}) a argument
  příkazu\note*{Marginální poznámka vysázená pomocí příkazu
    \pgm{\textbackslash note}.}  \pgm{\textbackslash note} bude zobrazen jako
  poznámka na okraji vysázená drobným bezpatkovým fontem.\footnote{Ve
    skutečnosti je marginální poznámka na této straně vysázena pomocí
    persistentní \uv{hvězdičkové formy} \pgm{\textbackslash note*}, která
    nebude ignorována ani při použití volby \opt{final}, což je pochopitelně
    případ i~tohoto dokumentu.}  Kromě toho je tato volba předána do všech
  balíčků,\INDEX{balíček} které ji implementují, takže např.
  balíček\INDEX{balíček} \pkg{graphicx}~\cite{Car:LIVE:CTAN:Graphics} zobrazí
  místo vložené grafiky obdélník příslušné velikosti s~uvedením názvu
  obrazového souboru včetně relativní cesty a podobně, čímž se kompilace a
  ladění urychlí.\footnote{Pokud navíc nechcete být během přípravy práce
    zdržováni sazbou titulních stránek, anotací, klíčových slov, poděkování,
    prohlášení, věnování, jednoduše zakomentujte příkazy \pgm{\textbackslash
      maketitle}, \pgm{\textbackslash tableofcontents}, \pgm{\textbackslash
      listoffigures}, \pgm{\textbackslash listoftables} a případně i~prostředí
    typu \pgm{annotation}, jež jsou součástí titulních stran. Příkazy
    \pgm{\textbackslash author}, \pgm{\textbackslash title} apod. mohou zůstat
    nezakomentovány, jejich obsah bude ignorován.}

\optitem{final}\label{opfinal} se používá u~kompilace finální verze práce;
  přetečené řádky nejsou označeny \emph{slimáky} (ostatně žádné by se neměly
  vyskytnout), a~argument příkazu \pgm{\textbackslash note} je
  \uv{anihilován}, takže marginálie vkládané pomocí \pgm{\textbackslash note}
  není třeba pro finální verzi odstraňovat. Balíček
  \pkg{graphicx}~\cite{Car:LIVE:CTAN:Graphics} vloží skutečné grafické
  soubory.

\optitem{htext}\label{opthtext} transformuje externí linky i~interní křížové
  odkazy na hypertext při zachování sazby nezávisle na nastavení \opt{draft}
  či \opt{final} nebo na kompilaci formátem \pkg{latex} (u~něj se předpokládá
  následná konverze do \designation{pdf} ovladačem
  \pgm{dvipdfmx}~\cite{Hir:LIVE:CTAN:dvipdfmx})\footnote{Chcete-li si výsledný
    \designation{dvi} soubor pouze prohlížet, volbu \opt{htext} vypněte.} či
  \pkg{pdflatex}.

\optitem{narrowmargin}\label{opnarrowmargin} rozšíří sazební obrazec o~cca
  $1.08\,\mathrm{in} \approx 27.4\,\mathrm{mm}$ a zvýší počet řádků na stránce
  z~39 na 41. Bez volby \opt{narrowmargin} je poměr šířky k~výšce 0.618 (zlatý
  řez), se zapnutou volbou \opt{narrowmargin} 0.707 (stříbrný řez).

\optitem{indentfirst}\label{opindentfirst} způsobí horizontální odsazování
  (indentaci) úvodního řádku prvního odstavce kapitoly, sekce, subsekce,
  \ldots\@ Implicitně se úvodní řádky prvního odstavce neindentují.

\optitem{fllastline}\label{optfllastline} potlačí centrování posledního
  (východového) řádku některých implicitně do bloku sázených odstavců. Typicky
  se to týká popisek obrázků a tabulek, anotace, klíčových slov, prohlášení,
  poděkování (nikoli věnování, které je celé sázeno centrovaně). Jejich
  východové řádky budou místo centrování zarovnány doleva \uv{flush left}.

\optitem{headfootsans}\label{optheadfootsans} změní font živého záhlaví a
  zápatí (na úvodních stránkách kapitol je v~něm umístěna stránková číslice
  místo v~záhlaví) ze skloněného patkového na skloněný bezpatkový font.
\end{description}

\noindent
Při kompilaci tohoto manuálu byly aktivní volby \opt{final} a \opt{htext}.

\endinput

%%
%% End of file `ca_docop.tex'.
